\section{相关工作}

\paragraph{大语言模型评测。}
GPT-3 的少样本学习结果表明，大规模语言模型可以通过纯文本提示完成广泛任务 \cite{brown2020language}。随后，MMLU 和 BIG-bench 等基准进一步将评测范围拓展至多学科知识、常识、数学和复杂任务集合 \cite{hendrycks2021measuring,srivastava2022beyond}。这些工作强调了跨任务能力测量的重要性，然而平均分数往往难以揭示模型在特定条件下的失效根因。本文选取一个边界更明确的空间推理场景，通过受控变量和重复采样刻画模型行为差异。

\paragraph{推理提示与行为测试。}
Chain-of-thought prompting 表明，中间推理步骤可以显著改变模型在复杂推理任务上的表现 \cite{wei2022chain}。然而可解释的推理文本并不必然意味着稳定的推理过程，尤其当题干中出现无关信息或近阈值参数时，模型可能生成看似合理但结论错误的解释。CheckList 提出从行为测试角度构造能力矩阵和测试类型，以避免仅凭总体准确率评测 \cite{ribeiro-etal-2020-beyond}；对抗阅读理解工作亦表明，无关或误导性文本可以显著破坏模型判断 \cite{jia-liang-2017-adversarial}。本文借鉴上述行为测试思想，将扰动、层级和反事实条件直接纳入 Prompt 设计。

\paragraph{空间推理。}
空间关系是人类认知和语言理解的基础。已有空间介词基准发现，大语言模型在空间关系推断上存在明显的提示敏感性和与人类表现的差距 \cite{comsa-narayanan-2023-benchmark}。与一般空间介词或视觉空间基准不同，本文聚焦一个纯文本、参数化、可判定的轮椅转弯任务：模型不需要识别图像，而需要在自然语言描述中提取几何约束，并判断轮椅能否完成转弯。这种任务适合观察语言模型在现实工程语境中的空间数量推理边界。
